% Complete author/root/B/D actual-file audits PASS; auxiliary object-chart input.
\subsection{Uniform generation and fixed-determinant charts on exceptional cycles}

\begin{theorem}[Proved: uniform generation of reflexive pullbacks]
\label{D:reflexive-pullback-uniform-generation}
Let $(A,\mathfrak m,k)$ be an excellent normal local surface
$k$-algebra with algebraically closed residue field $k$ and a
dualizing complex. Let $\pi:Y\to\operatorname{Spec}A$ be a proper
regular resolution, with exceptional curves $E_i$, and let $Z>0$
be an effective exceptional Cartier divisor of full support.
Assume that
\[
               -Z\cdot B\ge 2p_a(B)
       \quad\text{for every nonzero effective exceptional cycle }B.
                                                               \tag{D.UG.1}
\]
If $F$ is a vector bundle on $Y$ admitting a map
$\mathcal O_Y^a\to F$ which is surjective at every exceptional
generic point, then
\[
 H^1(Y,F(-Z))=0,\qquad
 H^1(Y,F(-Z)\otimes I_{P,Y})=0\quad(P\in E),             \tag{D.UG.2}
\]
and $F(-Z)$ is generated by its global sections. Here $P$ ranges
over closed points of the exceptional fibre.

In particular the assertion applies to
$F=(\pi^*M)^{\vee\vee}$ for every finite reflexive $A$-module
$M$, with a single choice of $Z$ independent of $M$ and its rank.
On a minimal resolution with
$\kappa_i=K_Y\cdot E_i\ge0$, it suffices to choose the cycle in
Lemma~\ref{D:integral-antinef-energy-cycle} so that
\[
 -Z\cdot E_i\ge3\kappa_i\quad(\kappa_i>0),
 \qquad -Z\cdot E_i\ge0\quad(\kappa_i=0).              \tag{D.UG.3}
\]
An entirely zero-canonical connected component can instead be
given its rational-double-point fundamental cycle.
\end{theorem}

\begin{proof}
\emph{The point-ideal form of automatic adjunction.}
Let $D$ be an effective exceptional cycle, $L$ a line bundle
on $D$, and $P$ a closed point of $D$. Suppose that
\[
                  \deg_B L\ge2p_a(B)
       \quad\text{for every nonzero subcycle }B\le D.    \tag{D.UG.4}
\]
Then
\[
       \operatorname{Hom}_D(I_{P,D}\otimes L,\omega_D)=0.
                                                               \tag{D.UG.5}
\]
We recall the proof, including the nilpotent layers. A nonzero
map on the left has a maximal scheme-theoretic supporting
subcurve $B$. Automatic adjunction factors it through
$\omega_B=\mathcal Hom_D(\mathcal O_B,\omega_D)$, with a
generically surjective map from its restricted source; see
\citet[Lemma~2.4 and the complete proof of Theorem~1.1,
pp.~192--193]{C_CFHR1999}. The annihilator has no embedded
associated point because the target is Cohen--Macaulay.
Since $Y$ is regular, a Cohen--Macaulay subcurve of its Cartier
divisor $D$ is precisely an effective Cartier subcycle $B\le D$.
The map from $I_{P,D}\otimes\mathcal O_B$ to
$I_{P\cap B,B}$ has a zero-dimensional kernel. That kernel
maps to zero in $\omega_B$. We therefore obtain a generically
surjective morphism
\[
                  I_{P\cap B,B}\otimes L|_B\longrightarrow\omega_B.
\]
Both sides have the same generic lengths; the source is
torsion-free. The morphism is consequently injective and has
zero-dimensional cokernel. Euler characteristics give
\[
 \deg_B L\le 2p_a(B)-2+\ell(P\cap B)\le2p_a(B)-1,
\]
contradicting (D.UG.4). The same argument with the empty cluster
gives $\operatorname{Hom}_D(L,\omega_D)=0$.
This is the full nonreduced-cycle argument, rather than an
inequality tested only on the reduced prime components.

\emph{Insertion of a generically generated vector bundle.}
Restrict the given evaluation map to $D=tZ$, where $t\ge1$.
It remains surjective at every generic Artinian stalk of $D$.
The target $\omega_D$ has no nonzero zero-dimensional subsheaf.
Precomposition therefore gives an injection
\[
 \operatorname{Hom}_D(I_{P,D}\otimes F|_D(-Z),\omega_D)
 \lhook\joinrel\longrightarrow
 \operatorname{Hom}_D(I_{P,D}\otimes\mathcal O_D(-Z),\omega_D)^a.
                                                               \tag{D.UG.6}
\]
Indeed a map in the kernel vanishes at every generic point,
so its image would be a zero-dimensional subsheaf of $\omega_D$.
By (D.UG.1), the right side vanishes by (D.UG.5). Serre duality
on the proper Gorenstein curve $D$ gives
\[
 H^1(D,F(-Z)|_D\otimes I_{P,D})=0,
 \qquad H^1(D,F(-Z)|_D)=0.                              \tag{D.UG.7}
\]

There is one Tor term to retain when passing from $Y$ to $D$:
\[
 \operatorname{Tor}_1^Y(k(P),\mathcal O_D)
       \longrightarrow I_{P,Y}\otimes\mathcal O_D
       \longrightarrow I_{P,D}\longrightarrow0.
                                                               \tag{D.UG.8}
\]
Its image is zero-dimensional. Tensoring with the vector
bundle $F(-Z)|_D$ shows that it does not change $H^1$.
Thus (D.UG.7) also computes the restrictions of
$F(-Z)\otimes I_{P,Y}$.
The ideals of $Z$ and of the closed fibre have the same
radical, so their powers are cofinal. The theorem on formal
functions now makes the completions of both finite $A$-modules
in (D.UG.2) zero. Faithfulness of completion, or Nakayama's
lemma, gives (D.UG.2) itself.

The point-ideal exact sequence implies surjectivity
$H^0(Y,F(-Z))\to F(-Z)\otimes k(P)$ at every $P\in E$.
The evaluation cokernel is coherent. Its proper image in the
local base, if nonempty, would contain the closed point, whereas
its support misses the entire closed fibre. The cokernel is
therefore zero. This proves global generation on $Y$.

For a reflexive $M$, take a finite generating set over $A$.
At a codimension-one regular stalk of $Y$, the natural map
$\pi^*M\to(\pi^*M)^{\vee\vee}$ is onto: the target is the
free quotient of a module over a discrete valuation ring.
Its generating sections therefore have the required generic
surjectivity. No claim that the original reflexive pullback
was already globally generated at every closed point is used.

Finally write $B=\sum m_iE_i$. Adjunction and negative
definiteness give
\[
          2p_a(B)=2+\kappa^{\mathsf t}m+B^2.
\]
If $\kappa^{\mathsf t}m\ge1$, then $B^2\le-1$ and the right
side is at most $\kappa^{\mathsf t}m+1\le2\kappa^{\mathsf t}m$.
Condition (D.UG.3) is stronger than (D.UG.1) in this case.
If $\kappa^{\mathsf t}m=0$, the support consists of smooth
rational $(-2)$-curves. Its integral intersection form is even
and negative definite, so $B^2\le-2$ and $p_a(B)\le0$.
Anti-nefness of $Z$ suffices. A lattice parameter $c\ge5$ in
Lemma~\ref{D:integral-antinef-energy-cycle} gives (D.UG.3).
For a whole zero-canonical component the same parity argument
applies with its fundamental cycle. This proves the last claim.
\end{proof}

\begin{proposition}[Proved: degree-linear fixed-determinant cycle charts]
\label{D:hereditary-fixed-determinant-cycle-chart}
Let $R$ be a fixed complete excellent normal local surface
$k$-algebra of characteristic $p>0$, with algebraically closed
coefficient and residue field $k$. Consider finite normal
quasi-etale cover domains $A/R$ of rank $n$ in the energy
family of Corollary~\ref{D:punctured-cover-linear-energy}.
Here the energy hypothesis means that the stated bounds for
$p_g$, the canonical energy and the total exceptional
self-intersection hold with constants depending only on
the fixed base. In particular one may complete the algebraic
local cover family in that corollary: excellence preserves
regularity of the base-changed resolution, and proper flat
base change preserves the exceptional curves, intersection
numbers and coherent-cohomology lengths. No energy assertion
for arbitrary nonalgebraizable complete covers is added here.
Write $Y$ for the minimal regular resolution of $A$,
and choose $Z$ as in Theorem~\ref{D:reflexive-pullback-uniform-generation}
and Corollary~\ref{D:punctured-cover-antinef-cycle}.
The choice depends on $A$, but not on a module or its rank.
Fix $t\ge1$ and put $D=tZ$.

Let $(M_a)_{a\in\mathbb Z/f\mathbb Z}$ be a cyclic system of
rank-$r$ reflexive $A$-modules with
$F_A^{[*]}M_a\simeq M_{a+1}$. Put
\[
       F_a=(\pi^*M_a)^{\vee\vee},\qquad V_a=F_a(-Z)|_D.
\]
There is a constant $c_t$, independent of $n,r,f$ and this
system, such that
\[
 0\le c_1(F_a)\cdot Z\le c_tnr,
 \qquad h^0(D,\det V_a)\le c_tnr.                    \tag{D.UG.9}
\]
Every $V_a$ is globally generated. For each fixed line bundle
$L=\det V_a$ on $D$, all globally generated rank-$r$ bundles
on $D$ with determinant isomorphic to $L$ occur as the
universal quotients over an open subset of
\[
                 H^0(D,L)^{\oplus(r+1)}.
                                                               \tag{D.UG.10}
\]
In particular the relevant fixed-determinant object chart has
at most $c_tnr(r+1)=O_t(nr^2)$ scalar coordinates. Its universal
quotient is an actual vector bundle, also over nonreduced
parameter schemes.

This is a statement about the objects on one finite exceptional
cycle. It asserts neither that every such bundle is the
restriction of a desired reflexive unit object, nor that a
unit object is determined by this restriction.
\end{proposition}

\begin{proof}
\emph{A fixed ideal and its energy.}
Use the nonzero fixed $\mathfrak m_R$-primary Noether-different
ideal $J$ of Proposition~\ref{D:cover-uniform-Noether-different}.
It factors multiplication on every reflexive $A$-module through
a finite free $A$-module. Let $C$ be the divisorial part of
$J\mathcal O_Y$. The line bundle $\mathcal O_Y(-C)$ is
generated at the exceptional generic points by the fixed
generators of $J$. Hence $C$ is anti-nef.

Principalize $J$ once on a fixed regular modification $Y_0$
of the base, writing $J\mathcal O_{Y_0}=\mathcal O_{Y_0}(-C_0)$.
Resolve the finite normalization above $Y_0$ and take a common
regular resolution $Y'$ dominating $Y$. The divisor of $J$
on $Y'$ is $C'=f^*C_0$, where $f:Y'\to Y_0$ has generic
degree $n$. Projection and the orthogonal decomposition for
the birational map $g:Y'\to Y$ give
\[
 (C')^2=nC_0^2,\qquad C'=g^*C+B_g,\qquad
 (C')^2=C^2+B_g^2\le C^2.
\]
The divisor $B_g$ is exceptional for $g$; no sign of its
coefficients is needed for the last inequality. Consequently
\[
                  -C^2\le n(-C_0^2).
                                                               \tag{D.UG.11}
\]
All these intersections are exceptional intersections over the
local base. The already proved energy estimate for $Z$ and
Cauchy--Schwarz for the negative intersection form imply
\[
       0\le\lambda:=-C\cdot Z
           \le\sqrt{(-C^2)(-Z^2)}=O(n).               \tag{D.UG.12}
\]

\emph{The determinant recurrence.}
The source factorization through a free module shows that
$J^{[p]}$ kills the ordinary-to-reflexive Frobenius cokernel.
Pull back the ordinary Frobenius linearization and extend
across codimension two on regular $Y$. Frobenius is flat on
$Y$, so the resulting inclusions of vector bundles are
\[
        F_{a+1}(-pC)\ \subset\ F_Y^*F_a\ \subset\ F_{a+1}.
                                                               \tag{D.UG.13}
\]
The first inclusion can be checked at every exceptional
discrete valuation ring: the least valuation of a $p$-th
power of a generator of $J$ is $p$ times the coefficient
of $C$. The resulting rational map is regular everywhere
by reflexivity. Thus the effective determinant divisor
$B_a$ of the second inclusion satisfies
\[
       c_1(F_{a+1})=p\,c_1(F_a)+B_a,
                    0\le B_a\le rpC.                 \tag{D.UG.14}
\]
Generic generation implies $\deg(\det F_a|_{E_i})\ge0$.
Put $x_a=c_1(F_a)\cdot Z\ge0$. Since $Z$ is anti-nef,
(D.UG.14) gives
\[
                   p x_a-rp\lambda\le x_{a+1}\le p x_a.
\]
Iterate around the $f$-cycle. Cancellation of $p^f-1$ yields
\[
                   0\le x_a\le\frac{rp\lambda}{p-1}=O(nr),
                                                               \tag{D.UG.15}
\]
independently of the period $f$.

\emph{Determinants on the full thickening.}
Theorem~\ref{D:reflexive-pullback-uniform-generation} gives
global generation of $F_a(-Z)$ and hence of $V_a$ and its
determinant. With $Q=-Z^2$, one has
\[
               \deg_D\det V_a=t(x_a+rQ)=O_t(nr).
\]
The full, possibly nonreduced cycle has
$h^0(\mathcal O_D)=O_t(n)$ by
Corollary~\ref{D:punctured-cover-antinef-cycle} and the
separate fundamental-cycle treatment of zero-canonical
components. Choose a global section of the globally
generated determinant avoiding the finitely many associated
components of the Cohen--Macaulay curve $D$. It is regular,
and its zero scheme is finite. The sequence
\[
     0\longrightarrow\mathcal O_D\longrightarrow\det V_a
                  \longrightarrow T_a\longrightarrow0
\]
has $\ell_k(T_a)=\deg_D\det V_a$. Therefore
\[
       h^0(\det V_a)\le h^0(\mathcal O_D)+\deg_D\det V_a
                         =O_t(nr).
\]
This retains all nilpotent layers; it is not merely a degree
calculation on the reduced fibre.

\emph{The actual object chart.}
Let $V$ be globally generated of rank $r$ on $D$ and fix an
isomorphism $\det V\simeq L$. A general $r+1$-tuple of global
sections generates $V$: the rank-drop condition has codimension
two at a fixed point, while the reduced support is one-dimensional.
The incidence argument avoids every point of that support,
and Nakayama gives the same conclusion on $D$ with its
nilpotents. The kernel of the surjection
$\mathcal O_D^{r+1}\twoheadrightarrow V$ is a line bundle,
necessarily isomorphic to $L^{-1}$. Conversely a tuple of
$r+1$ sections of $L$ generating $L$ defines a subbundle
\[
        L^{-1}\longrightarrow\mathcal O_D^{r+1}
\]
whose quotient is locally free of rank $r$ and has determinant
$L$. The generating condition is open: its complement is the
proper image of the common zero scheme in the product with $D$.
The universal quotient is locally a split quotient of vector
bundles, so it remains a vector bundle after every base change.
The dimension bound follows from (D.UG.9).
\end{proof}

\begin{remark}[The exact remaining realization gate]
The charts in (D.UG.10) are for a fixed determinant line, not
for a moving Picard scheme. They need no scheme of rank-$r$
quotients with a quadratic rank enlargement. Nevertheless,
their open parameter spaces contain bundles which need not
come from a reflexive Frobenius-periodic module. A proof of
the hereditary representation bound still needs an effective
finite-thickening realization and isomorphism criterion,
including all Frobenius arrows and Artinian families, with
the same degree-linear complexity. Reduced-branch framing
does not by itself settle geometric normal thickenings or
the full formal overlap. No representation count or new
finite-generation conclusion is asserted here.
\end{remark}
